\bookchapter{Chapitre 4 : Installer et gérer les applications — le « pack de compétences » du smartphone}

\chapterintro{Les applications rendent le smartphone capable de tout faire, mais il faut aussi bien les gérer. Comme vous assigneriez un poste et des autorisations à un nouvel employé, l'installation d'applications nécessite un processus standardisé.}

\booksection{Boutique officielle : le « supermarché d'applications » le plus sûr}

\begin{itemize}
  \item \textbf{Android}: Google Play (à l'étranger) ou les boutiques intégrées de chaque marque (comme Huawei AppGallery, Xiaomi App Store).
  \item \textbf{iOS}: Uniquement l'App Store.
  \item \textbf{HarmonyOS}: Huawei AppGallery.
\end{itemize}

Importance sécuritaire essentielle : examinées, mises à jour garanties, paiements sécurisés. Guide d'action :\textbf{téléchargez en priorité uniquement depuis les boutiques officielles !}

\booksection{Mode pur/protection : activez le « pare-feu antivirus » du smartphone}

Les grandes marques nationales (Huawei, Xiaomi, OPPO, vivo, Honor) intègrent toutes un\textbf{mode pur}ou\textbf{mode protection}, un commutateur de sécurité au niveau système.\textbf{Il est fortement recommandé à tous les utilisateurs de l'activer}.

\booksubsection{Qu'est-ce que c'est exactement ?}

La fonction principale du mode pur (aussi appelé « mode protection ») est de\textbf{bloquer l'installation d'applications tierces}: lorsqu'il est activé, le système\textbf{n'autorise l'installation d'applications que depuis la boutique officielle}, les packages d'installation téléchargés depuis les navigateurs, QR codes, applications de messagerie, etc. sont directement bloqués, éliminant les logiciels malveillants à la source.

\booksubsection{Public concerné}

\begin{itemize}
  \item \textbf{Utilisateurs seniors}: Incapables de distinguer les publicités « officielles » des « contrefaçons », le mode pur est la protection la plus fiable.
  \item \textbf{Utilisateurs peu sensibles aux autorisations d'applications}: N'ont pas besoin d'installer des applications de niche, versions bêta, jeux étrangers.
  \item \textbf{Tous les utilisateurs souhaitant améliorer la sécurité de leur appareil}: Même les passionnés, l'activer vaut mieux que ne pas l'activer — en cas de besoin temporaire, il suffit de le désactiver ponctuellement.
\end{itemize}

\booksubsection{Chemin d'activation (exemple Huawei/Xiaomi)}

\begin{itemize}
  \item \textbf{Huawei/HarmonyOS}: [Paramètres] > [Système et mise à jour] >\textbf{[Mode pur]}, sélectionnez « Activer ».
  \item \textbf{Xiaomi}: [Paramètres] > [Protection de la vie privée] >\textbf{[Protéger la vie privée]}ou [Paramètres] > [Paramètres des applications] > [Installation d'applications], désactivez « Installer des applications inconnues ».
  \item \textbf{OPPO/vivo}: [Paramètres] > [Sécurité] >\textbf{[Installer des applications de sources externes]}, désactivez les autorisations une par une.
\end{itemize}

\booksubsection{Mécanisme de contrôle sécurité de la boutique d'applications}

La boutique d'applications n'est pas là pour la décoration, elle dispose d'une chaîne de sécurité complète :

\begin{itemize}
  \item \textbf{Analyse antivirus}: Chaque application est soumise avant publication à une analyse conjointe de plusieurs moteurs antivirus pour identifier les chevaux de Troie et virus connus.
  \item \textbf{Vérification des autorisations}: La boutique examine si les autorisations demandées par l'application sont « raisonnables » — une application de lampe de poche qui demande l'accès aux contacts ? Rejet ou retrait immédiat.
  \item \textbf{Certification du développeur}: La publication nécessite une certification d'entreprise, de carte d'identité personnelle ou de carte de crédit. En cas de problème, on peut retrouver la personne, ce ne sont pas des « applications fantômes ».
  \item \textbf{Surveillance continue}: Les applications déjà publiées sont régulièrement réexaminées, tout comportement malveillant nouvellement découvert entraîne un retrait forcé.
\end{itemize}

\booksection{Sources tierces et « parasites électroniques » : risque extrêmement élevé !}

Télécharger des applications depuis des pages web ou des liens inconnus peut inclure des virus, voler vos données privées, frauder vos paiements. Mentionnons ici un nouveau concept : les\textbf{« parasites électroniques »}.

\textbf{Qu'est-ce qu'un « parasite électronique » ?}\par

Il s'agit de\textbf{logiciels « crackés » implantés de chevaux de Troie}— vous croyez « économiser », en réalité vous « tombez dans le piège ». Formes courantes :

\begin{itemize}
  \item Versions « crackées », « gratuites », « sans publicité » de logiciels payants ;
  \item Outils gris comme « boost de likes », « capture d'enveloppes rouges », « tricheurs de jeux » ;
  \item Prétendues versions « internes » d'applications recommandées par des « petites annonces » dans des groupes de discussion.
\end{itemize}

Ces programmes semblent fonctionner, mais en coulisses ils peuvent\textbf{voler les contacts, lire les SMS, écouter le micro, miner en arrière-plan}. Le mode pur est la\textbf{première ligne de défense}contre eux — en vous empêchant directement de les installer.

\booksubsection{Deuxième ligne de défense recommandée :}

\begin{itemize}
  \item \textbf{Analyse antivirus complète régulière}: Utilisez le centre de sécurité intégré au smartphone (comme Huawei Gestionnaire de téléphone, Xiaomi Centre de sécurité) pour une analyse mensuelle.
  \item \textbf{Accord prudent des autorisations sensibles}: Contacts, SMS, localisation, appareil photo, micro — ces cinq « autorisations à haut risque » doivent être strictement contrôlées.
  \item \textbf{Surveillez les anomalies de consommation et de données}: Le smartphone chauffe soudainement, la batterie se vide vite, les données consommées la nuit — il y a probablement un « parasite » caché en arrière-plan.
\end{itemize}

Si vous devez absolument installer une application tierce :

\begin{enumerate}
  \item Assurez-vous que la source est absolument fiable.
  \item Lisez attentivement les exigences d'autorisation avant l'installation.
  \item Analysez avec une application de sécurité après l'installation.
\end{enumerate}

Conseil ultime : dans les paramètres système, désactivez l'autorisation « Installer des applications inconnues ».

\booksection{Installer, mettre à jour, désinstaller}

\begin{itemize}
  \item \textbf{Installation}: Cliquez sur « Obtenir » ou « Installer » dans la boutique, à la première ouverture après installation, les autorisations vous seront demandées.
  \item \textbf{Mise à jour}:\textbf{Il faut mettre à jour !}Principalement pour corriger les failles de sécurité. Il est recommandé d'activer les mises à jour automatiques.
  \item \textbf{Désinstallation}: Appuyez longuement sur l'icône et glissez vers la zone de désinstallation. Pour un nettoyage plus complet, allez dans [Paramètres] > [Gestion des applications] pour trouver et désinstaller.
\end{itemize}

\booksection{Gestion des autorisations : vous êtes le « patron » des applications}

Les applications demandent des autorisations (comme contacts, localisation) parfois à juste titre (une carte a besoin de la localisation), parfois sans raison (une lampe de poche a besoin des contacts).

\begin{itemize}
  \item \textbf{Options d'autorisation}: « Toujours autoriser », « Autoriser uniquement pendant l'utilisation », « Interdire », « Demander à chaque fois ».
  \item \textbf{Principes de gestion}:
  \begin{enumerate}
    \item Minimisation : ne pas accorder ce qui n'est pas nécessaire.
    \item Audit a posteriori : vérifiez régulièrement dans les paramètres d'autorisation qui « dépasse ses droits ».
    \item Contextualisation : comme pour une application VTC, n'accordez la localisation que pendant l'utilisation.
  \end{enumerate}
\end{itemize}

\begin{tcbwarn}{Autorisations à surveiller en priorité}
\begin{itemize}
  \item Contacts/SMS : peuvent être utilisés pour envoyer des spams.
  \item Localisation en temps réel : vos déplacements peuvent être suivis.
  \item Appareil photo/Micro : risque d'enregistrement à votre insu.
  \item Accès aux fichiers : vos fichiers privés peuvent être lus.
\end{itemize}

\end{tcbwarn}

\booksection{Éléments constitutifs d'un package d'installation}

Vous êtes-vous déjà demandé ce que contient ce package d'installation téléchargé depuis la boutique d'applications ? Comprendre la constitution du package d'installation vous aide à mieux comprendre comment les applications fonctionnent, et à savoir où chercher en cas de problème.

\booksubsection{Structure d'un fichier APK Android}

Un APK (Android Package Kit) est essentiellement une archive ZIP. Après décompression, on trouve les répertoires et fichiers principaux suivants :

\begin{itemize}
  \item \textbf{AndroidManifest.xml}: La « carte d'identité » et le « mode d'emploi » de l'application. Il enregistre le nom du package, le numéro de version, les autorisations requises, la version système prise en charge, l'Activity point d'entrée et d'autres informations essentielles. Le système lit ce fichier en premier lors de l'installation.
  \item \textbf{classes.dex}: Le code principal de l'application. Le code source Java/Kotlin compilé devient du bytecode Dalvik, stocké dans ce fichier. Une application peut avoir plusieurs fichiers dex (comme classes2.dex, classes3.dex).
  \item \textbf{resources.arsc}: La table d'index des ressources compilées. Enregistre les ID de toutes les ressources et leurs correspondances.
  \item \textbf{Répertoire res/}: Dossier de ressources. Contient images, fichiers de mise en page, chaînes de caractères, définitions de couleurs, etc. Les images de différentes résolutions sont placées dans différents dossiers drawable (comme drawable-hdpi, drawable-xhdpi).
  \item \textbf{Répertoire assets/}: Dossier de ressources brutes. Stocke les ressources non compilées, comme fichiers web, polices, audio, etc. L'application peut les lire via AssetManager à l'exécution.
  \item \textbf{Répertoire lib/}: Dossier de bibliothèques natives. Stocke les fichiers .so compilés (produits de compilation C/C++), classés par architecture CPU (comme arm64-v8a, armeabi-v7a). Les jeux et les applications nécessitant des calculs haute performance ont généralement ce répertoire.
  \item \textbf{Répertoire META-INF/}: Dossier d'informations de signature. Contient les trois fichiers MANIFEST.MF, CERT.SF, CERT.RSA, utilisés pour vérifier l'intégrité et la signature de l'APK.
\end{itemize}

\booksubsection{Structure d'un fichier IPA iOS}

L'IPA (iOS App Store Package) est également une archive compressée, avec une structure légèrement différente :

\begin{itemize}
  \item \textbf{Répertoire Payload/}: Où se trouve le contenu principal. Il contient un dossier .app qui inclut tout le contenu de l'application.
  \item \textbf{Info.plist}: Fichier de configuration de l'application, similaire à l'AndroidManifest.xml d'Android. Enregistre l'ID de l'application, la version, les autorisations requises, les appareils pris en charge, etc.
  \item \textbf{Fichier exécutable}: Programme principal du même nom que le dossier .app, un fichier binaire au format Mach-O compilé.
  \item \textbf{Répertoire Frameworks/}: Bibliothèques dynamiques et frameworks dont dépend l'application.
  \item \textbf{Fichiers de ressources}: Images, audio, vidéo et autres ressources, généralement stockées dans le dossier .app.
  \item \textbf{Répertoire \_CodeSignature/}: Informations de signature de code, utilisées pour vérifier l'intégrité et l'origine de l'application.
\end{itemize}

\booksubsection{Mécanisme de vérification de signature}

Pourquoi les applications téléchargées depuis la boutique officielle sont-elles plus sûres ? La vérification de signature est la clé.

\begin{itemize}
  \item \textbf{Signature Android}: Le développeur signe l'APK avec une clé privée, le système vérifie avec la clé publique lors de l'installation. Si l'APK est altéré (par exemple avec un virus injecté), la signature devient invalide et le système refuse l'installation. Google Play exige également App Signing, où Google conserve la clé de signature, protégeant davantage la sécurité de l'application.
  \item \textbf{Signature iOS}: L'examen d'Apple est encore plus strict. Chaque IPA doit avoir une signature avec certificat délivré par Apple, sinon l'installation est impossible. C'est pourquoi iOS ne permet d'installer des applications que depuis l'App Store ou via certificat d'entreprise. Les applications signées avec un certificat d'entreprise ont une durée de validité limitée, et Apple peut les révoquer à tout moment.
  \item \textbf{Rôle de la signature}:
  \begin{enumerate}
    \item Vérification de l'origine : confirmer que l'application provient bien du développeur déclaré.
    \item Vérification de l'intégrité : s'assurer que l'application n'a pas été altérée pendant le transfert.
    \item Prévention de l'usurpation : une application malveillante ne peut se faire passer pour une autre.
  \end{enumerate}
\end{itemize}

\booksubsection{Taille du package d'installation et optimisation}

Vous avez peut-être remarqué que la taille de l'application affichée dans la boutique et l'espace occupé après installation ne sont pas les mêmes. C'est parce que :

\begin{itemize}
  \item \textbf{Taille de téléchargement}: Taille de l'APK/IPA après compression.
  \item \textbf{Taille d'installation}: Taille après décompression, généralement bien supérieure à la taille de téléchargement.
  \item \textbf{Occupation à l'exécution}: L'application génère également des caches et fichiers de données à l'exécution, occupant davantage d'espace.
\end{itemize}

Moyens d'optimisation des applications modernes :

\begin{itemize}
  \item \textbf{App Bundle (Android)}: Technologie de distribution dynamique de Google Play. L'utilisateur ne télécharge que les ressources adaptées à son appareil (par exemple uniquement les images correspondant à la résolution de son écran), réduisant considérablement la taille de téléchargement.
  \item \textbf{App Thinning (iOS)}: Technologie similaire, Apple élimine automatiquement les ressources inutiles selon l'appareil.
  \item \textbf{Téléchargement à la demande}: Les applications volumineuses comme les jeux peuvent d'abord télécharger le package principal, puis les ressources supplémentaires ensuite.
\end{itemize}

\begin{tcbinfo}{Le saviez-vous : pourquoi certaines applications sont-elles si volumineuses ?}
Les applications de jeu sont généralement les plus volumineuses, car elles contiennent de nombreuses textures HD, modèles 3D et fichiers audio. Les applications sociales sont également assez lourdes, car elles doivent prendre en charge divers modules fonctionnels. Les applications utilitaires sont généralement les plus légères. Si vous remarquez qu'une application utilitaire est anormalement volumineuse, méfiez-vous de fonctionnalités inutiles intégrées.

\end{tcbinfo}

\begin{tcbpractice}{Exercice pratique : examiner les autorisations de 3 applications}
Choisissez 3 applications courantes (comme sociale, shopping, utilitaire), accédez à la page de gestion des autorisations des applications dans [Paramètres], examinez chaque autorisation et désactivez ou modifiez en « Autoriser uniquement pendant l'utilisation » les autorisations sensibles non nécessaires.

\end{tcbpractice}
